\chapter{Аудит происхождения коэффициентов кривизны}
\label{ch:v10-cell-birth-curvature-coefficient-origin-audit}

\section{Восемь путей происхождения}

Предыдущий гейт показал, что кривизненный родитель выбирает
$q=\ell^2$ только через отношение $B/(2A)$. Проверим восемь кандидатов на
происхождение пары коэффициентов: космологическую кривизну, спектральное
обрезание, энергию часов, обратную длину ячейки, собственную кривизну,
температуру КМС, топологическую плотность и индуцированную петлю материи.

Каждый кандидат проверяется по шести условиям: правильные парные
размерности, внутреннее происхождение, независимость от искомой длины,
вывод общим родителем, типизированное отображение в $(A,B)$ и разрыв
абсолютной масштабной орбиты. Матрица кандидатов имеет ранг $3$, но вектор
полных проходов равен
\begin{equation}
 (0,0,0,0,0,0,0,0)^T.
 \label{eq:v10-curvature-coefficient-pass-vector}
\end{equation}
Наибольший результат $4/6$ получают спектральное обрезание и петля
физической материи. Однако первый путь вводит свободное обрезание, а второй
ещё не выведен существующим общим родителем.

\section{Общий запрет масштабно-ковариантной пары}

Пусть кандидат сначала поставляет величину $m$ размерности $L^{-2}$, а
затем строит
\begin{equation}
 A=\alpha m^2,
 \qquad B=\beta m,
 \label{eq:v10-curvature-coefficient-seed}
\end{equation}
где $\alpha$ и $\beta$ безразмерны. Тогда
\begin{equation}
 \frac{B^2}{A}=\frac{\beta^2}{\alpha},
 \qquad
 q_*=\frac{\beta}{2\alpha m},
 \qquad
 q_*m=\frac{\beta}{2\alpha}.
 \label{eq:v10-curvature-coefficient-invariant}
\end{equation}
Родитель выбирает лишь безразмерное произведение $qm$. Преобразование
\begin{equation}
 q\mapsto s^2q,
 \qquad m\mapsto\frac{m}{s^2}
 \label{eq:v10-curvature-seed-orbit}
\end{equation}
сохраняет весь функционал. Эквивалентная карта для $(A,B,q)$ имеет
ранг $2$, ядро размерности $1$ и образующий $(-2,-1,1)^T$.

\begin{theorem}[Размерно правильная пара ещё не является происхождением]
Любой кандидат вида $(A,B)=(\alpha m^2,\beta m)$ воспроизводит правильные
масштабные степени и условный минимум, но не выбирает абсолютную величину
$m$. Он переносит исходную масштабную свободу в коэффициенты.
\end{theorem}

\begin{nogocriterion}[Запрет кругового коэффициентного якоря]
Кривизну, энергию часов, температуру или обратную длину нельзя использовать
для вывода $A$ и $B$, если сама эта величина уже вычислена через искомый
масштаб ячейки. Полный проход требует независимого физического механизма,
создающего коэффициент и нарушающего орбиту.
\end{nogocriterion}

\begin{protocol}[Статус аудита коэффициентов]\begin{enumerate}
\item проверенных кандидатов: \textbf{$8$};
\item полных проходов: \textbf{$0/8$};
\item ранг матрицы кандидатов: \textbf{$3$};
\item наибольший результат: \textbf{$4/6$};
\item происхождение пары $(A,B)$: \textbf{$0/1$};
\item абсолютный масштаб ячейки: \textbf{$0/1$};
\item следующий гейт: \textbf{родитель индуцированных гравитационных коэффициентов физической материи}.
\end{enumerate}\end{protocol}

Машинный сертификат сохранён в
\path{s2t_v10_cell_birth_four_volume_curvature_coefficient_origin_candidate_audit_gate_results.json}.